\documentclass[11pt]{article}
\usepackage[utf8]{inputenc}
\usepackage{amsmath,amssymb}
\usepackage{enumitem}
\usepackage{geometry}
\geometry{margin=1in}

% Custom macros for questions
\newcommand{\QBlock}[3]{%
  \subsection*{Question #1}
  #2
  \begin{enumerate}[label=(\Alph*)]
  #3
  \end{enumerate}
  \vspace{1em}
}

\newcommand{\Correct}[2]{%
  \item \textbf{#1} \hfill $\checkmark$ \\
  \textit{Explanation: #2}
}

\newcommand{\Wrong}[2]{%
  \item #1 \\
  \textit{Explanation: #2}
}

\begin{document}

\title{Isomorphic Physics Questions}
\author{Generated by AI}
\date{\today}
\maketitle

% Questions go here using \QBlock, \Correct, \Wrong macros
\QBlock{1}{Which equation represents the correct conversion?}{
  \Wrong{1\,¢ = 1\,¥}{A cent is not equivalent to one yen.}
  \Wrong{1 = 0.1}{This is a mathematically incorrect statement.}
  \Correct{1\,¢ = 0.007\,¥}{1\,¢ is approximately 0.007 yen, making this conversion accurate.}
  \Wrong{100\,¢ = 100\,¥}{100 cents do not equal 100 yen.}
}

\QBlock{2}{Identify the incorrect unit for the quantity listed.}{
  \Wrong{Work: electron volt}{Electron volt is a valid unit for work.}
  \Wrong{Power: watt}{Watt is the correct unit for power.}
  \Correct{Momentum: meter/second}{The correct unit for momentum is kilogram\cdot meter/second.}
  \Wrong{Charge: coulomb}{Coulomb is the correct SI unit for charge.}
}

\QBlock{3}{A bird flies from the Wall to a tree located at position $\vec{r} = 4\hat{i} + 3\hat{j}$ kilofurlongs. Evaluate the distance between the Wall and the tree.}{
  \Correct{The distance is 5 kilofurlongs.}{Distance is found using the Pythagorean theorem: $\sqrt{4^2 + 3^2} = 5$ kilofurlongs.}
  \Wrong{The distance is 7 kilofurlongs.}{The calculation error arises from incorrect root evaluation.}
  \Wrong{The distance is 3 kilofurlongs.}{Misused components result in an underestimated distance.}
  \Wrong{None of the above is correct.}{Option A is accurate as per the Pythagorean theorem.}
}

\QBlock{4}{An object's displacement is given by $s(t) = (2.1 \frac{m}{s^4})t^4 + (3.5 \frac{m}{s^3})t^3$. Which statement is true about acceleration?}{
  \Wrong{This object has constant acceleration.}{The presence of a $t^4$ term indicates a non-constant acceleration.}
  \Wrong{The magnitude of the acceleration is 3.5 $m/s^3$.}{Acceleration is time-dependent, determined through further differentiation.}
  \Correct{The acceleration at $t = 0$ is 0 $m/s^2$.}{Acceleration, calculated through double differentiation, is zero at $t = 0$.}
  \Wrong{The acceleration at $t = 0$ is 4.2 $m/s^2$.}{Misinterpretation of initial condition leads to incorrect acceleration value.}
}

\QBlock{5}{An object moves between position A $(\vec{r} = -6\hat{i} + 8\hat{j})$ and B $(\vec{r} = 1\hat{i} - 2\hat{j})$ over 4 seconds. Determine the correct statement about the journey.}{
  \Correct{The moving distance is 12.5 m.}{The distance is calculated as $\sqrt{(1 - (-6))^2 + (-2 - 8)^2} = 12.5$ meters.}
  \Wrong{The speed at $t = 0$ must be 3.1 m/s.}{Incorrect speed as it does not reflect the average journey speed.}
  \Wrong{The acceleration at $t = 0$ must be 1.0 $m/s^2$.}{Non-applicable as initial acceleration data is not provided.}
  \Wrong{None of the above is correct.}{Option A accurately describes the distance.}
}

\QBlock{6}{A bird takes off at a speed $v_0$ and angle $\theta$ to the horizontal. What is the speed at its peak height?}{
  \Wrong{$v_0$.}{Complete speed remains reduced to horizontal component at peak altitude.}
  \Correct{$v_0\cos\theta$.}{The vertical velocity is zero, leaving $v_0\cos\theta$ as the speed at peak.}
  \Wrong{$v_0\sin\theta$.}{Vertical component becomes zero at maximum height.}
  \Wrong{0.}{Horizontal velocity persists, preventing zero speed.}
}

\QBlock{7}{A drone traveling at $(25 \frac{m}{s})\hat{i}$ is provided with an acceleration $(3 \frac{m}{s^2})\hat{i} + (5 \frac{m}{s^2})\hat{j}$. Which statement is accurate about its motion?}{
  \Wrong{The x-component of the drone's velocity remains constant.}{Acceleration impacts x-component, causing it to change.}
  \Wrong{The y-component of the drone's velocity remains constant.}{Y-velocity is affected due to the y-acceleration.}
  \Correct{Both speed and direction alter.}{Acceleration modifies both the speed and the motion's path.}
  \Wrong{None of the above is correct.}{Option C is the complete and correct description.}
}

\QBlock{8}{If each charge's magnitude is doubled while their distance is also doubled, what happens to the electrostatic force?}{
  \Correct{It remains unchanged.}{The increased charge and distance cancel each other out in the formula $F \propto \frac{q_1q_2}{r^2}$.}
  \Wrong{It decreases by a factor of $\sqrt{2}$.}{This option ignores balanced effects on force.}
  \Wrong{It increases by a factor of 2.}{This misstates balance between charge and distance increments.}
  \Wrong{It increases by a factor of 4.}{Exaggerated force increase disregards square distance effect.}
}

\QBlock{9}{An object escapes a planet with radius R and mass M at escape velocity $v = \sqrt{2GM/R}$. Which statement holds true?}{
  \Wrong{A greater mass object requires higher escape speed.}{Escape speed relies on planetary properties, not the escaping object.}
  \Wrong{A lesser mass object requires lower escape speed.}{Unrelated to the object's mass; it depends on planetary attributes.}
  \Correct{Larger mass and smaller radius planets require greater escape speed.}{A planet's mass increases, and radius decreases, requiring more speed.}
  \Wrong{Lesser mass and larger radius planets need greater escape speeds.}{This configuration reduces required escape speed.}
}

\QBlock{10}{A force $\vec{F} = (5.02\, \text{N})\hat{i} + (2.04\, \text{N})\hat{j}$ acts over $x = 0.5\,\text{m} \text{ to } x = 6.5\,\text{m}$. Calculate the work done.}{
  \Correct{30.1 J}{Work is calculated along x-axis, $(5.02)(6.0) = 30.1\,\text{J}$.}
  \Wrong{15.3 J}{Misinterpretation of force-action span.}
  \Wrong{7.00 J}{Underestimation in applied work.}
  \Wrong{42.0 J}{Inflated figure calculated through misapplied force duration.}
}

\QBlock{11}{A 6 kg ball initially traveling at 2.5 m/s hits a wall and rebounds at 1.5 m/s backward. The impact lasts for 0.3 seconds. Determine the average force exerted on the ball.}{
  \Correct{60 N}{Change of momentum $(m(v_f - v_i))/t$ yields 60 N.}
  \Wrong{30 N}{Incorrect partial use of velocity change within force calculations.}
  \Wrong{75 N}{Enhanced momentum change miscalculations.}
  \Wrong{None of the above is correct.}{Option A reflects correct force application.}
}

\QBlock{12}{A goalie, originally at rest, catches a 200 g puck initially moving at 160 km/hr. What is the goalie's speed after catching the puck if his mass is 80 kg?}{
  \Correct{0.089 m/s}{Linear momentum conservation gives $m_{puck}v_{puck} = (m_{goalie}+m_{puck})v_{goalie}$.}
  \Wrong{0.072 m/s}{This results from incorrect simultaneous speed adjustment.}
  \Wrong{120 m/s}{Typographical error in perceived post-collision speed result.}
  \Wrong{200 m/s}{Interface misuse leads to high errant speed calculation.}
}

\end{document}